\documentclass[conference]{IEEEtran}
\usepackage[utf8]{inputenc}
\usepackage[T1]{fontenc}
\usepackage{cite}
\usepackage{amsmath,amssymb}
\usepackage{graphicx}
\usepackage{booktabs}
\usepackage{url}

\title{Can Deterministic Network Verification Reduce Undetected Faults in LLM‑Generated Configurations?}

\author{
\IEEEauthorblockN{Autonomous AI Researcher}
\IEEEauthorblockA{CNSM 2026 Agentic AI Researcher Track}
}

\begin{document}

\maketitle

\begin{abstract}
We preregistered and executed a paired confirmatory study (study id c1-gnr-vv-2026-0001) to measure whether inserting a deterministic network verifier between an LLM intent\ensuremath{\rightarrow}configuration generator and an emulated execution reduces the proportion of configurations that would execute with operational faults. Using shared\_initial\_candidate semantics (one identical LLM candidate per intent evaluated both without and with verification), we evaluated 72 preregistered paired intents with gpt-5-mini and a canonical deterministic verifier (deterministic\_netops\_validator\_v5). Execution completed with 144 episodes (72 pairs) and 72 model calls. All baseline executions produced no emulator-observed faults (n\_11 = 72, n\_10 = n\_01 = n\_00 = 0). The paired difference in undetected-fault proportion is 0.0 (paired bootstrap 95\% CI [0.0, 0.0], exact McNemar two-sided p = 1.0; bootstrap seed = 1729, resamples = 10000). Under the preregistered shared\_initial\_candidate contract this degenerate confirmatory outcome is expected when identical generated candidates produce no baseline execution faults. We report the preregistration rationale, deterministic execution accounting, representative validator diagnostics, exact inferential outputs, and artifact-grounded recommendations for follow-on confirmatory designs able to detect verifier utility under realistic baseline-error regimes.
\end{abstract}

\section{Introduction}
Generative large language models (LLMs) are increasingly proposed to translate high-level operator intents into device-level network configurations. This intent\ensuremath{\rightarrow}configuration automation can speed change pipelines but risks silently violating reachability, access-control, ordering, or workflow invariants and thereby causing outages if generated text is executed without additional checks. A standard operational mitigation is a generate--verify--repair (GVR) architecture that combines stochastic LLM generation with deterministic verification and, if needed, repair or human review. However, despite many architectural proposals and prototypes, systematic preregistered experiments that measure a verifier's detection performance against executed outcomes under tightly controlled paired designs are scarce.

This paper answers a narrowly scoped, preregistered research question: does inserting a deterministic network verifier between an LLM generator and emulator execution materially reduce the proportion of configurations that execute with operational faults when the identical generated candidate is evaluated both without and with verification? That estimand isolates verifier detection from improvements that could arise from re-sampling, prompt-engineering, or repair-enabled iterations. We executed study c1-gnr-vv-2026-0001 using gpt-5-mini and deterministic\_netops\_validator\_v5 across 72 preregistered paired intents under shared\_initial\_candidate semantics. The run completed with no technical failures, produced a degenerate confirmatory outcome (no baseline emulator faults), and yields a paired difference of 0.0 (95\% CI [0.0,0.0], exact McNemar p = 1.0). Rather than treating this as a null failure, we analyze why it occurred given the preregistered contract and archived artifacts, and we provide concrete, artifact-grounded guidance for designs that can detect verifier utility when baseline error prevalence is non-negligible.

\section{Related Work}
Work on LLM-driven NetOps spans intent\ensuremath{\rightarrow}configuration translation, multi-agent/tool-augmented orchestration, and architectural proposals for combining stochastic generation with deterministic checks. Prior empirical studies demonstrate that modern LLMs can produce device-level configurations from high-level intent in constrained vendor-dialect settings and curated evaluation suites; these results show promise but are typically scoped to limited task families and curated datasets rather than broad production traces [1]. That line of work motivates leveraging LLMs to reduce operator effort while highlighting the need for safety controls when configuration semantics and ordering matter.

Complementing single-agent generation work, multi-agent and tool‑augmented frameworks have been proposed that explicitly integrate verifiers, tool calls, or specialized modules into LLM-driven NetOps pipelines to reduce regressions and coordinate multi-step changes [2]. These systems illustrate practical architectures for inserting deterministic checks or domain services into agentic workflows, and their evaluations commonly rely on emulators, curated scenarios, or prototype deployments. Such prototypes suggest feasible integration points for verifiers but do not, in the supplied records, provide preregistered paired measurements of verifier detection performance against executed faults.

A related and more general strand of work articulates and evaluates generate--verify--repair (GVR) patterns outside NetOps---particularly in code- and software-repair domains---where stochastic generators produce candidates that deterministic checkers validate and, when necessary, trigger repair iterations [3]. These studies document how deterministic validation can convert nondeterministic generation into a correctness-convergent workflow; they also show that the empirical utility of verification depends strongly on the baseline error rate of generator outputs and on whether the architecture permits verifier-triggered repairs or resampling. Translating these insights into NetOps requires coupling LLM outputs to domain-appropriate verifiers and measuring detection relative to executed outcomes.

Canonical deterministic network-verification techniques---exemplified by header-space analysis and related formal reachability/policy analyzers---provide precise, semantics-aware invariants for reachability and ACL-style policy checks and are natural candidates for the 'verify' stage in GVR architectures [4]. Such verifiers can deterministically flag violations of specified invariants or workflow constraints expressed over packet headers, interface states, and rule orderings. However, the literature lacks, among the supplied records, preregistered paired experiments that measure how often these canonical verifiers would have detected faults that would otherwise manifest when LLM-generated configurations are executed in an emulator or live device.

Taken together, the verified literature establishes three pertinent points: (i) LLMs can produce plausible device-level configs from intent in constrained settings [1]; (ii) system proposals and prototypes show how verifiers and tools can be integrated into multi-agent NetOps workflows [2]; and (iii) GVR patterns can yield empirical correctness improvements in other domains when baseline generator errors are non-negligible and when verifier-triggered repairs or resampling are allowed [3]. What remains underexplored in the supplied evidence is a preregistered, paired measurement that isolates verifier detection on identical candidates---i.e., does a canonical deterministic verifier reduce the probability of an execution fault for the same LLM-generated candidate when the candidate is executed with and without prior verification? Our study directly targets this measurement gap by adopting a shared\_initial\_candidate paired design and reporting exact execution-accounting artifacts to make that detection question reproducible and auditable.

\section{Methodology}
Preregistration and primary estimand. We preregistered and froze study c1-gnr-vv-2026-0001 before observing outcomes. The primary estimand is the paired reduction in undetected operational-fault proportion (paired\_success\_rate\_difference\_guarded\_minus\_baseline). The confirmatory hypothesis (H1) targeted an absolute paired reduction of 0.20 (two-sided alpha = 0.05, power >= 0.80); a sample-size heuristic yielded \ensuremath{\approx}63 pairs and we preregistered N = 72 pairs (24 per topology-difficulty stratum). The preregistration and analysis plans, including sensitivity analyses and exclusion rules, are part of the archived preregistration record.

Shared\_initial\_candidate semantics and experimental contract. The execution\_contract enforced shared\_initial\_candidate semantics: for each intent we performed exactly one initial LLM generation and evaluated that identical candidate in two paired conditions. Baseline condition: execute the shared candidate in an isolated emulator and record whether emulator-observed operational faults occur. Guarded condition: run the canonical deterministic verifier on the same candidate; if the verifier flags violations the guarded outcome is recorded as prevented (no executed fault); if the verifier does not flag, execute the identical candidate in the emulator and record the result. By construction, any observed paired difference can only arise from the verifier preventing execution of a candidate that would otherwise have failed at emulation; this isolates detector utility from generator-side resampling or repair effects.

Model, adapter, and verifier configuration (artifact-grounded). The declared model is gpt-5-mini (execution/model\_configuration.json), requested via the hosted\_netops\_gvr\_v1 adapter. Key recorded model parameters: declared\_model\_version = "gpt-5-mini", provider = "openai\_responses", max\_output\_tokens = 2000, maximum\_attempts\_per\_call = 1, and temperature = null. The recorded temperature = null indicates that provider-default runtime sampling semantics were used; we treat this as an archived sampling-parameter uncertainty rather than an explicit numeric temperature. The canonical verifier is deterministic\_netops\_validator\_v5 (validator\_version = "5.0") configured to run the full\_intent\_constraint\_validation rule set. The verifier emits structured validator\_traces per episode (fields include parsed\_command\_count, required\_command\_count, normalized\_configuration, validation\_before and validation\_after final\_state snapshots, violation\_codes, violation\_count, and difficulty metadata). The verifier's validity verdict and trace for each episode were archived under execution/scoring/ and formed the mechanistic basis for guarded decisions.

Task-bank construction, contamination controls, and stratification. Confirmatory tasks were generated deterministically by deterministic\_netops\_task\_generator\_v5. Each task\_manifest entry encodes: a natural-language intent, a machine-structured required\_sequence (ordered field changes that implement the intent), allowed\_touched\_fields, preserved-state policies, initial and expected\_final\_state snapshots, and a difficulty.level tag (medium/hard/very\_hard). The preregistration required deterministic exact-match contamination checks against a public-corpus fingerprint (analysis/contamination\_summary.csv); exact-match flagged items are excluded from confirmatory analysis (none were excluded in this run). The confirmatory sample followed the preregistered stratification (24 pairs per difficulty stratum) to allow descriptive subgroup inspection.

Execution protocol, retries, and deterministic accounting. The missingness\_plan permitted up to two retries (three attempts total) for transient hosted-API errors on generation calls; all retries and provider events were logged. The run started at 2026-08-31T06:56:23.065707Z and completed at 2026-08-31T07:06:28.665518Z (execution manifest). The execution manifest records planned\_episode\_count = 144, completed\_episode\_count = 144, failed\_episode\_count = 0, and model\_calls\_used = 72 (one shared generation per pair). Provider-call audit (deterministic\_reconciliation.json) reports hosted\_model\_call\_event\_count = 72, completed\_hosted\_model\_call\_event\_count = 72, failed\_hosted\_model\_call\_event\_count = 0, cache\_miss\_count = 72, and stage\_counts = \{"shared\_initial\_generation": 72\}. These deterministic accounting artifacts were used to reconcile paired outcomes and to confirm that the shared\_initial generation stage was the only cross-condition generation activity.

Scoring, validation rules, and per-episode traces. For each episode the verifier produced a binary validity verdict (valid / invalid) and an accompanying validator\_trace JSON that documents parsed and required command counts, normalized\_configuration, validation\_before/after states, and any violation\_codes. The scoring rule deployed was full\_intent\_constraint\_validation (validator enforces required\_sequence order, allowed\_touched\_fields, and preserve\_unspecified\_state constraints encoded in the task manifest). Per-episode scoring artifacts (execution/scoring/task-*.json) therefore permit audit of why a candidate was accepted or flagged. In the recorded run no verifier flagging occurred for confirmatory pairs (violation\_count = 0 across representative episodes), and no automated repair was applied (repair\_applied = false in validator\_trace fields).

Inferential and sensitivity analysis procedures (deterministic scripts). The primary analysis used the registered paired\_binary\_analysis\_v1 executor to construct the 2\ensuremath{\times}2 contingency table (guarded \ensuremath{\times} baseline) using binary indicators where an undetected executed fault is 1 and success/no-fault is 0. The prespecified exact inferential test is an exact McNemar two-sided test when discordant pairs exist. We also computed paired bootstrap-percentile 95\% confidence intervals for the paired difference using 10,000 resamples (bootstrap\_seed = 1729), recorded in analysis/analysis\_log.jsonl. Pre-specified subgroup confirmatory comparisons (topology complexity and policy type) were registered with Holm correction to control familywise error; these controls are implemented in the deterministic analysis pipeline but become non-informative in the absence of discordant pairs. Pre-registered sensitivity analyses executed by the same scripts include (a) treating excluded confirmatory pairs as failures (complete\_pair\_primary\_with\_failure\_as\_zero\_sensitivity) and (b) bootstrap diagnostics for detector detection and false-positive rates. Because the run produced no exclusions and no discordant outcomes these sensitivity analyses reproduce the degenerate numerical conclusion.

Design-level notes and construct tradeoffs (artifact-linked). The shared\_initial\_candidate semantics intentionally isolates verifier detection from generator-side sampling; this preserves a clean causal interpretation of verifier-prevention utility but prevents measurements that rely on verifier-triggered resampling or repair. The archived execution/model\_configuration.json recording temperature = null documents that provider-default sampling behavior was allowed; had we required explicit sampling parameters (non-null temperatures) we could have measured sampling-sensitivity as a separate factor. The verifier's full\_intent\_constraint\_validation rule set enforces task-manifest constraints: because the task generator encoded explicit required\_sequence and allowed\_touched\_fields many generated candidates closely matched those constraints, which the verifier therefore did not flag. This alignment is measurable in validator\_trace fields (parsed\_command\_count equals required\_command\_count and violation\_count = 0 for representative episodes) and explains the ceiling result when baseline emulator faults are absent. The preregistered sensitivity and alternative contracts (independent-generation arms, repair-enabled flows, adversarial task seeding, and multi-model comparisons) are described in the preregistration and can be executed using the same deterministic execution and analysis infrastructure archived with this run.

\section{Results}
Execution and manifest summary. The autonomous pipeline completed without technical failures. Key manifest figures (execution/ and analysis/ manifests): planned\_episode\_count = 144; completed\_episode\_count = 144; failed\_episode\_count = 0; model\_calls\_used = 72; artifact\_count = 510. analysis/condition\_summary.csv reports baseline successes = 72, baseline failures = 0 (success\_rate = 1.0) and guarded successes = 72, guarded failures = 0 (success\_rate = 1.0).

Paired contingency and exact inference. Aggregating across the 72 confirmatory pairs produced contingency counts n\_11 = 72, n\_10 = 0, n\_01 = 0, n\_00 = 0 (guarded \ensuremath{\times} baseline). The paired estimate paired\_success\_rate\_difference\_guarded\_minus\_baseline = 0.0. Exact McNemar two-sided p = 1.0. The paired bootstrap percentile 95\% CI (seed = 1729; 10,000 resamples) is [0.0, 0.0], reflecting the degenerate distribution produced by the absence of discordant pairs. analysis/paired\_contingency\_table.csv and analysis/condition\_summary.csv contain these tabulations and are archived.

Representative validator diagnostics and per-episode exemplars. Representative per-episode scoring artifacts are archived (execution/scoring/task-000001-*.json ... task-000006-*.json). Representative summaries (extracted from archived JSONs):
- pair-000001 (task-000001): difficulty.level = "medium", parsed\_command\_count = 4, required\_command\_count = 4, normalized\_configuration matched the shared candidate, validator\_version = "5.0", violation\_count = 0.
- pair-000002 (task-000002): difficulty.level = "hard", parsed\_command\_count = 2, required\_command\_count = 2, violation\_count = 0.
- pair-000003 (task-000003): difficulty.level = "hard", parsed\_command\_count = 6, required\_command\_count = 6, violation\_count = 0.
These traces show the verifier executed on each shared candidate and produced no violation flags across the sampled representative tasks. Execution/scoring/*.json files contain full validator\_traces for reproducibility.

Contamination, missingness, and sensitivity diagnostics. analysis/contamination\_summary.csv reports zero exact-match contaminations for confirmatory pairs. analysis/missingness\_summary.csv records complete\_pairs = 72 and zero failed or unscorable episodes. Pre-specified sensitivity analyses (treating excluded pairs as failures; bootstrap diagnostics) were executed and archived; because no pairs were excluded and no discordant outcomes occurred these sensitivity analyses reproduce the degenerate numeric conclusion (paired difference = 0.0). analysis/analysis\_log.jsonl records bootstrap\_seed = 1729 and bootstrap\_resamples = 10000.

\section{Discussion}
Confirmatory interpretation and boundary conditions. The preregistered paired design with shared\_initial\_candidate semantics validly measures verifier detection on identical candidates; because the identical generated candidates produced zero emulator faults the verified paired outcome is a ceiling/null result: the verifier could not reduce executed faults that did not occur. The numeric outputs (n\_11 = 72; paired difference = 0.0; McNemar p = 1.0; bootstrap CI [0.0,0.0]) follow directly from the preregistered contract and observed data. Reporting this degenerate outcome is important: it documents a reproducible case where the preregistered estimand and data-generating conditions leave no room for the verifier to demonstrate utility.

Artifact-grounded diagnostic factors. Archived artifacts allow concrete attribution of this ceiling outcome to measurable pipeline factors. First, the deterministic task generator (deterministic\_netops\_task\_generator\_v5) produced structured manifests with explicit required\_sequence and allowed\_touched\_fields that constrain acceptable configs; these manifest constraints appear in execution/task\_manifest.jsonl and increase the probability of a single correct generation. Second, the declared model (gpt-5-mini recorded in execution/model\_configuration.json) produced candidates that consistently matched those structured constraints across the synthetic corpus (execution/responses/* files and per-episode shared\_initial\_candidate\_sha256 entries). Third, the verifier's rule set (full\_intent\_constraint\_validation in deterministic\_netops\_validator\_v5) enforces the same required\_sequence and field constraints encoded in the manifests; per-episode validator\_traces under execution/scoring/ show parsed\_command\_count = required\_command\_count and violation\_count = 0 for representative tasks. Together these aligned design choices---constrained intents, congruent verifier rules, and a model that produced matching sequences---explain the absence of baseline emulator faults.

Statistical and design implications. From a statistical perspective, a paired confirmatory test that targets a fixed absolute reduction (here preregistered 0.20) requires a non-zero baseline prevalence of executed faults to have power to detect a reduction; when baseline prevalence is zero the paired difference must be zero regardless of verifier performance. This is not a paradoxical failure of the verifier or analysis but an expected algebraic consequence of the estimand under shared\_initial\_candidate semantics. Practically, therefore, confirmatory designs intended to measure verifier detection should ensure the task bank yields a measurable baseline error rate (e.g., by including ambiguous intents, adversarial cases, or known failure-mode seeds) or should permit condition-specific sampling or repair so that the verifier can alter the executed candidate distribution.

Concrete preregisterable alternatives to detect verifier utility. The archived pipeline supports several clear modifications---each compatible with preregistration---that would permit measuring different aspects of verifier utility without inventing new artifacts here: (1) independent-generation arms: allow separate initial generator calls per condition (baseline vs guarded) so that the verifier (or guarded workflow) can influence candidate distributions via repairs or sampling differences; (2) repair-enabled flows: permit a single deterministic repair call when the verifier flags a violation and then execute the repaired candidate to measure end-to-end improvement in execution outcomes; (3) adversarial/fault-injection task banks: seed tasks with intentionally ambiguous wording, ordering subtleties, or vendor-edge cases to raise baseline error prevalence and probe verifier detection under stress; (4) sampling-parameter sweeps and multi-model comparisons: record explicit non-null temperature and test multiple model versions to quantify model- and sampling-sensitivity; and (5) dedicated false-positive cost arms: measure operational blocking cost by comparing outcomes when verifier blocks flagged candidates versus when blocked candidates are repaired and re-executed. All of these are operationally meaningful, preregistrable changes and can be implemented using the same evidence-recording conventions shown in our run (execution/model\_configuration.json, execution/scoring/, analysis/).

Threats to validity revisited with evidence. Internal validity: shared\_initial\_candidate semantics intentionally isolates verifier detection but prevents verifier-mediated resampling or repair effects; the deterministic provider-call audit (deterministic\_reconciliation.json) confirms one shared generation per pair, making the isolation explicit. Construct validity: emulator-executed faults are a conservative surrogate for operational harm but do not capture traffic-dependent timing or vendor-hardware idiosyncrasies---limitations recorded in the manuscript. External validity: the run used a single declared model (gpt-5-mini) and a synthetic task bank; extrapolation to other models, richer real-world task distributions, or production hardware remains limited. Conclusion validity: with zero baseline faults the preregistered power analysis targeted to an absolute 0.20 reduction becomes moot; future confirmatory designs must align expected baseline prevalence with power targets.

Operational implications and measured tradeoffs. The observed ceiling does not imply that verifiers are unnecessary in practice. Instead, it shows that under constrained, high-clarity intents and a verifier rule set that encodes the same constraints, an LLM can produce device-level sequences that pass deterministic checks and emulator execution. In operational settings where intents are incomplete, telemetry is noisy, or vendor dialects vary, verifiers remain a principled guard. Measuring the verifier's net operational value requires explicit, preregistered trade-off quantification: estimate true-positive detection rates on a task bank with non-negligible baseline faults, and quantify verifier false-positive blocking costs (the proportion of flagged-but-safe candidates and operational delay or repair effort). The archived artifacts and deterministic accounting in this study provide a reproducible template to design such targeted, powered follow-on experiments.

\section{Limitations}
\begin{itemize}
\item Confirmatory ceiling/null outcome: the baseline generator produced zero emulator faults on the preregistered task set, preventing direct measurement of verifier detection benefit.
\item Shared\_initial\_candidate semantics: the design measures verifier effect on the same generated candidate only and does not assess independent per-condition generation, multi-sample generation, or repairs unless explicitly permitted by the contract.
\item Single-model scope: experiments used one declared model (gpt-5-mini); results do not imply generality across models or versions.
\item Synthetic/emulated tasks: task corpus and emulator executions differ from production devices and live traffic; external validity to production deployments is limited.
\item Sampling-parameter uncertainty: execution/model\_configuration.json shows temperature = null (sampling parameters unset); null indicates provider-default runtime sampling semantics and is a residual uncertainty recorded in the archived configuration.
\item Residual contamination risk: deterministic provenance and exact-match checks were applied, but private training-data contamination of the hosted model cannot be fully ruled out and remains residual uncertainty.
\end{itemize}

\section{Conclusion}
We executed a preregistered paired evaluation (c1-gnr-vv-2026-0001) under shared\_initial\_candidate semantics to test whether a canonical deterministic verifier reduces the proportion of LLM-generated configurations that would execute with operational faults. For the chosen model (gpt-5-mini), verifier (deterministic\_netops\_validator\_v5), and synthetic/emulated task bank, baseline executions produced zero emulator faults and the verifier did not change outcomes (paired difference = 0.0; bootstrap 95\% CI [0.0,0.0]; exact McNemar p = 1.0). This artifact-grounded null/ceiling outcome is internally consistent with the preregistered design: when identical generated candidates produce no executed faults a verifier cannot reduce executed faults under the shared\_initial\_candidate contract. Measuring verifier utility under realistic error regimes requires preregistered contracts that permit independent or multiple generator samples, repair-enabled flows, adversarial task sets, and multi-model comparisons. The archived artifacts (responses, task manifests, validator traces, execution manifest, and analysis logs) provide a reproducible baseline and a template for those follow-on confirmatory designs and power calculations.

\section*{Disclosure Statement}
Autonomy and artifacts: This study was executed autonomously after an autonomy lock. Human role: a Research Director selected the topic family and approved the immutable initial master prompt prior to the autonomy lock; no human manually edited technical manuscript content after the lock. Models and tools: initial generation used gpt-5-mini (declared\_model\_version recorded in execution/model\_configuration.json) orchestrated via the hosted\_netops\_gvr\_v1 adapter; deterministic\_netops\_validator\_v5 (validator\_version = "5.0") served as the canonical verifier; an isolated emulator executed candidate configurations. Preregistration: study c1-gnr-vv-2026-0001 was preregistered and frozen before observing outcomes. Immutable master prompt reference: provenance/master\_prompt.txt (SHA-256: 1872df1e\allowbreak{}1805d2d9\allowbreak{}6940456c\allowbreak{}a016bd66\allowbreak{}5d1d5196\allowbreak{}add77f5a\allowbreak{}cdf1582b\allowbreak{}b39b15ba). Key archived artifacts include execution/raw\_results.jsonl, execution/task\_manifest.jsonl, execution/model\_configuration.json, analysis/paired\_contingency\_table.csv, analysis/condition\_summary.csv, deterministic\_reconciliation.json, and per-episode validator traces under execution/scoring/. The model configuration records temperature = null (provider-default sampling semantics archived in execution/model\_configuration.json). Provider-call audit: hosted\_model\_call\_event\_count = 72. No human manual manuscript editing or content insertion occurred after the autonomy lock.

\begin{thebibliography}{99}
\bibitem{ref1} Nguyen Tu, Sukhyun Nam, James Won‐Ki Hong, "Intent‐Based Network Configuration Using Large Language Models", 2024, 10.1002/nem.2313
\bibitem{ref2} Zhaodong Wang, Samuel Lin, Guanqing Yan, Soudeh Ghorbani, Minlan Yu, Jiawei Zhou, Nathan Hu, Lopa Baruah, Sam Peters, Srikanth Kamath, Jian Yang, Ying Zhang, "Intent-Driven Network Management with Multi-Agent LLMs: The Confucius Framework", 2025, 10.1145/3718958.3750537
\bibitem{ref3} Rafael Cadenas, "Convergent Correctness in Stochastic Code Generation: A Generate-Verify-Repair Architecture for Deterministic Validation of Autonomous AI Coding Agents in Regulated Environments", 2026, 10.2139/ssrn.6754899
\bibitem{ref4} http://citeseerx.ist.psu.edu/viewdoc/summary?doi=10.1.1.300.8659
\end{thebibliography}

\end{document}
